\documentclass[11pt,a4paper]{article}
\usepackage[margin=2.4cm]{geometry}
\usepackage{amsmath,amssymb,booktabs,hyperref}
\newcommand{\Lc}{\textbf{(c)}\ }
\newcommand{\Ld}{\textbf{(d)}\ }
\newcommand{\scr}{\texttt{audit\_\allowbreak toda\_\allowbreak dictionary.py}}
\newcommand{\dd}{\mathrm d}
\emergencystretch=2em
\title{An exact endpoint diagnostic for the Toda momentum and phase dictionary}
\author{Quantum Boussinesq solitons project}
\date{4 November 2026}
\begin{document}
\maketitle
\begin{abstract}
The implemented Toda Fredholm equations imply an exact edge momentum and an
endpoint reciprocity identity for their directional phases. \Lc\ These fix
the symmetric part of a possible quantum frame correction. They do not fix
its antisymmetric part or identify it with a physical charge transformation.
\Ld\ An independent 60-digit calculation finds a nonseparable symmetric
remainder after the classical frame subtraction actually used in the
overtaking code. \Lc\ This shows that the retained classical subtraction does not
antisymmetrise the directional kernel at finite $\hbar$; it does not
by itself refute the forward ordered physical amplitude or determine the
edge limit. The existing model, momentum prescription, and normalizations are
retained; no replacement physical convention is selected. \Ld
\end{abstract}

\section{Definitions and scope}
All results below are consequences of the existing asymptotic Bethe-ansatz
integral equations, not a new proof of their physical applicability. \Lc
denotes project derivations and checks; \Ld\ marks unresolved physical
identifications. The core high-precision checker \scr\ is independent;
an isolated final comparison uses the original implementation.
It records input hashes in \texttt{audit\_\allowbreak toda\_\allowbreak dictionary\_\allowbreak results.json}.
No external literature claim is introduced.

Use $h=\hbar_T$, $t=2\pi h$, and the quantum ground-state interval $[-q,q]$.
The code \texttt{toda\_\allowbreak edge\_\allowbreak ba.py:26--67} defines
\begin{align}
 \Theta(d)&=2h\,\operatorname{Im}\log\Gamma(1+id/h)+2d\log h,
 &K(d)&=\Theta'(d),\\
 (\mathcal A f)(x)&=tf(x)-\int_{-q}^q K(x-y)f(y)\,\dd y,
 &\mathcal A\rho&=1,\qquad \int_{-q}^q\rho=1/a,\\
 \mathcal A F_p(x)&=\Theta(p-x),
 &P(p)&=p+\int_{-q}^q\Theta(p-x)\rho(x)\,\dd x,\\
 \Phi(p,z)&=\Theta(p-z)-\int_{-q}^q K(p-x)F_z(x)\,\dd x.
 \label{eq:definitions}
\end{align}
Here $p$ is the bare Bethe spectral parameter. $P(p)$ is its dressed
momentum. Assume the ground-state integral operator is invertible and its
solutions are smooth up to the finite endpoints, as in the tested case.
The even kernel makes $\mathcal A$ self-adjoint for the unweighted integral
inner product. The ground density is even, and $P$ is odd.

\section{The edge subtraction in momentum is exact}
Differentiation and the ground-state equation give, inside the sea,
$P'(p)=t\rho(p)$. Integrating this identity from $-q$ to $q$ and using
oddness yields \Lc\ (\scr)
\begin{equation}
 \boxed{P(q)=\frac{\pi h}{a},\qquad P(-q)=-\frac{\pi h}{a}.}
 \label{eq:edgeP}
\end{equation}
Self-adjointness also gives
$P(p)=p+\int_{-q}^q F_p(x)\,\dd x$, the total momentum including backflow.
Thus the relative physical momentum used by the code is a subtraction of
the known quantum Fermi momentum, rather than a new definition of the bare
coordinate. \Lc

For a right-moving excitation, the implemented label is fixed through
\begin{equation}
 P(p)-P(q)=\frac{4b}{a}(\kappa\cosh\kappa-\sinh\kappa),\quad
 b=e^{-a/2},\quad \kappa=\sqrt{yh/b},\quad u=\frac{p-q}{h}.
 \label{eq:dictionary}
\end{equation}
These are the existing $\kappa=\sqrt3\beta$,
$h=3b\hbar_B$, $y=\beta^2/\hbar_B$ conventions. Differentiating at fixed
ground state gives an exact Jacobian \Lc\ (\scr)
\begin{equation}
 \frac{\dd u}{\dd y}=\frac{2\sinh\kappa}{aP'(p)}.
 \label{eq:Jacobian}
\end{equation}
This specifies what an asymptotic derivation of $u(y)$ must invert;
identifying $u$ with dressed physical momentum would skip this Jacobian.
Equation~\eqref{eq:edgeP} alone neither proves nor disproves the proposed
large-$y$ logarithmic expansion for $u(y)$. \Ld

There is a small implementation distinction: \texttt{edge()} returns
$e=q(1+10^{-10})$, not $q$. Its displacement in edge units is exactly
$(e-q)/h=10^{-10}q/h$. \Lc\ The offset is tiny in the saved runs but is
not uniform in a literal $h\to0$ limit. A controlled edge derivation must
specify that numerical offset tends to zero faster than $h$, or use the
exact endpoint analytically. This audit changes no production code. \Ld

\section{Exact reciprocity and its proof}
For arbitrary real external arguments define the symmetric dressed kernel
\[
 \mathcal R(p,z)=K(p-z)+\int_{-q}^q
       K(p-x)\,[\mathcal A^{-1}K(z-\cdot)](x)\,\dd x.
\]
Differentiate \eqref{eq:definitions}, and integrate once by parts in the
equation for $F_z$. The resulting identities are \Lc\ (\scr\ checks the
integrated consequences below)
\begin{align}
 \partial_z\Phi(p,z)&=-\mathcal R(p,z),\\
 (\partial_p+\partial_z)\Phi(p,z)
   &=\mathcal R(p,q)F_z(q)-\mathcal R(p,-q)F_z(-q),\\
 \partial_pF_p(\pm q)&=\mathcal R(\pm q,p)/t.
\end{align}
Consequently the derivative in either argument of
$\Phi(p,z)+\Phi(z,p)$ equals that of the following endpoint bilinear.
The integration constant vanishes at $p=z=0$ by parity. Hence \Lc
\begin{equation}
 \boxed{\Phi(p,z)+\Phi(z,p)
       =t\{F_p(q)F_z(q)-F_p(-q)F_z(-q)\}.}
 \label{eq:reciprocity}
\end{equation}
This retains the finite-sea boundary terms. Discarding them would incorrectly
make the directional Bethe phase antisymmetric.

Let $e_1,e_2$ be either exact edge, and set
\begin{align}
 \widehat\Phi(p,z)&=\Phi(p,z)-\Phi(p,e_2)-\Phi(e_1,z)+\Phi(e_1,e_2),\\
 f_p^\pm&=F_p(\pm q)-F_{e_1}(\pm q),\qquad
 g_z^\pm=F_z(\pm q)-F_{e_2}(\pm q).
\end{align}
The reverse phase uses exchanged references. Subtracting four copies of
\eqref{eq:reciprocity} gives \Lc\ (same-side and opposite-side checks)
\begin{equation}
 \widehat\Phi(p,z)+\widehat\Phi(z,p)
       =t(f_p^+g_z^+-f_p^-g_z^-).
 \label{eq:four}
\end{equation}
For clarity, write this as $\mathfrak M_1\mathfrak Q_2+
\mathfrak M_2\mathfrak Q_1$ using the computable response functions
\begin{align}
 \mathfrak M(p;e)&=ta\rho(q)\{F_p(q)-F_p(-q)-F_e(q)+F_e(-q)\},\\
 \mathfrak Q(p;e)&=\frac{F_p(q)+F_p(-q)-F_e(q)-F_e(-q)}{2a\rho(q)}.
 \label{eq:responses}
\end{align}
Integration of the differentiated $F$ equation gives
$t\rho(q)[F_p(q)-F_p(-q)]=1-P'(p)$, so equivalently \Lc
\begin{align}
 \mathfrak M(p;e)&=-a[P'(p)-P'(e)].
 \label{eq:Mresponse}
\end{align}
These symbols denote spectral response functions, not newly established
physical compression and momentum charges. \Ld\ Their normalization is
useful for comparison: the classical expressions have
$P_0'(p)=2\kappa/a$, $P_0'(2b)=0$, and
$-a\partial_aP_0(p)|_p=P_0(p)-2b\sinh\kappa$, verified symbolically in
\scr. \Lc\ This does not prove a uniform limit of the quantum endpoint
functions.

Any frame correction intended to make the four-term phase antisymmetric
must have the symmetric part \eqref{eq:four}. The product
$\mathfrak M_1\mathfrak Q_2$ is an algebraic example with that symmetric
part. Adding an arbitrary antisymmetric two-body function leaves the
condition unchanged. No such product is selected here as the physical
quantum frame, and the identity supplies no Ward identity relating these
functions to continuum quantum charges. \Ld

\section{The actual two-term overtaking subtraction fails this test}
The overtaking code \texttt{toda\_\allowbreak edge\_\allowbreak twoterm.py:11--17} uses only
$N(p,z)=\Phi(p,z)-\Phi(p,e_2)$. For its reverse directional expression,
an exact rearrangement gives \Lc\ (\scr)
\begin{align}
 N(p,z)+N(z,p)
 &=\widehat\Phi(p,z)+\widehat\Phi(z,p)+U(z)+V(p),\\
 U(z)&=\Phi(e_1,z)-\Phi(e_1,e_2),\qquad
 V(p)=\Phi(e_2,p)-\Phi(e_2,e_1).
\end{align}
Thus two-term and four-term normalizations have identical mixed dependence;
they are not the same phase. A rectangular mixed difference removes $U,V$
and any other one-body reference without selecting a convention.

At each fixed dressed momentum, retain the original classical frame
\[
 C_{\rm cl}(\kappa_1,\kappa_2)
 =(-2\kappa_1)\left[\frac{P_n(\kappa_2)}a-2b\sinh\kappa_2\right].
\]
Let $S_{\rm rem}=N(p_1,p_2)+N(p_2,p_1)-C_{\rm cl}(1,2)-C_{\rm cl}(2,1)$,
where the bare $p_j$ are obtained by inverting \eqref{eq:dictionary}.
For $\kappa_1\in\{0.8,1.0\}$ and $\kappa_2\in\{0.35,0.5\}$, define
the mixed rectangle as the upper--upper minus lower--upper minus
upper--lower plus lower--lower values. At $h=0.4$, $q=1.5$, the count
equation fixes $a=1.160312923180467\ldots$. Independent calculations give
\Lc
\begin{center}
\begin{tabular}{rrr}\toprule
Gauss nodes & Largest endpoint identity error & Rectangle of $S_{\rm rem}/h$\\\midrule
40 & $1.93\,10^{-10}$ & $-0.1125712665447065$\\
64 & $4.42\,10^{-16}$ & $-0.1125712665272591$\\
88 & $1.09\,10^{-21}$ & $-0.1125712665272590$\\\bottomrule
\end{tabular}
\end{center}
The original \texttt{EdgeBA} implementation at 120 nodes, including its
fixed edge offset, gives $-0.1125712665368\ldots$, differing by less than
$10^{-11}$. The same independent rectangle is obtained directly from
$\mathfrak M_1\mathfrak Q_2+\mathfrak M_2\mathfrak Q_1
-C_{\rm cl}(1,2)-C_{\rm cl}(2,1)$.
At the fixed spectral-to-dressed-momentum dictionary used here, it cannot
be removed by one-body phase/reference subtractions, including switching
between the displayed two- and four-term phase normalizations. \Lc
Changing the momentum reference and re-inverting the bare parameters is a
different operation and is not tested by this rectangle.
The reverse directional phase here is used as a reciprocity diagnostic;
it is not asserted to be a second overtaking collision in the same
ordering sector. This test is necessary for a scalar antisymmetric
frame identification, not a complete derivation of that identification.
\Ld

Any extension to a common antisymmetric scalar frame therefore requires
an additional mixed symmetric correction at finite $h$. It does not determine whether its normalized correction
vanishes, remains finite, or diverges in the particular nested edge and
large-$Y$ limits. The direct Toda expansion should use the implemented
quantity as defined; comparison to continuum physical scattering still
requires the antisymmetric frame part and the physical charge dictionary.
\Ld

\section{Verification and failures recorded}
\scr\ uses exact SymPy identities and a fresh 60-digit mpmath Nystr\"om
solver. Fixing $(h,q)$ and determining $a$ from the count condition selects
the same ground-state family as fixing $(h,a)$ and solving for $q$.
The dense quadrature, rather than decimal arithmetic, limits the displayed
accuracy. Checks include backflow momentum, edge momentum, reciprocity,
both reference orientations, the response derivative, physical-momentum
inversion, and the actual classical-frame mixed rectangle. Intermediate
resolvent and endpoint derivative identities are also checked directly. All pass.

\begin{center}\small
\begin{tabular}{p{.40\linewidth}p{.07\linewidth}p{.40\linewidth}}
\toprule Claim & Label & Verification\\\midrule
Dressed edge momentum and backflow & (c) & Independent Fredholm solve\\
Momentum Jacobian and classical responses & (c) & Exact symbolic derivatives\\
Reciprocity and endpoint bilinear & (c) & Mesh convergence and direct derivative checks\\
Two-/four-term phase relation & (c) & Exact rearrangement and numerical comparison\\
Classical-frame mixed remainder & (c) & Three meshes and original-code comparison\\
Physical frame and nested edge limit & (d) & Not established by these identities\\
\bottomrule
\end{tabular}
\end{center}

Failures recorded:
\begin{enumerate}
\item Exploratory 24- and 40-node meshes gave reciprocity discrepancies
about $3.5\,10^{-7}$ and $6.3\,10^{-11}$. Mesh refinement was required;
the identity and subtraction were not changed to suppress those errors.
\item The implemented classical frame fails the nonseparable symmetric
two-body test above. This failure is retained as the substantive result.
\item An initial PDF layout edit mismatched an equation environment. The
LaTeX build stopped; the environment was corrected before the final build.
\item No implementation assertion failed, and no test tolerance was
weakened. The fixed relative edge offset is recorded explicitly rather
than treated as an exact endpoint in a formal limit.
\end{enumerate}
The script ends with \texttt{ALL CHECKS PASSED} for the stated identities
and detected discrepancy. Physical interpretation of the remaining
antisymmetric frame and of the continuum quantum charges remains open. \Ld
\end{document}
